\begin{abstract}
Photographing a canteen tray and receiving a rough energy range is a practical
way to raise nutritional awareness among students, but a canteen serves a menu
that no public dataset contains: the dishes are Chinese, the photographs are
hand-held, and a twenty-way classifier always answers, even when the plate holds
something it has never seen. This paper reports a controlled study of that
setting. We start from CNFOOD-241, a 241-class Chinese dish dataset, and
pre-train a ResNet-18 on \cnPreTrainN{} images before fitting a 21-way canteen
head (twenty dishes plus a catch-all bucket) on \cnTrainN{} images. Pre-training
raises test accuracy from \cnNoPreTest\% to \cnTest\% with the same recipe, and a
six-view test-time fusion adds \cnFuseGain{} points on top, reaching
\cnFuseTop\% top-1 and \cnFuseTopFive\% top-5 on \cnTestN{} held-out photographs.
The twenty dish rows of the confusion matrix reach \cnDishAcc\%, but the
catch-all bucket is only \cnOtherRecall\% recalled, which is why we treat
open-set behaviour as a first-class result rather than a footnote: a domain
probe on pooled features plus a Chinese-dish veto decides whether the system
names a dish at all. On five evaluation sets the gate keeps \gateIn\% of canteen
photographs, \gateVal\% of the validation split, but admits only \gateOutCn\% of
fifteen unseen Chinese dishes, \gateOutOther\% of non-food scenes and
\gateReal\% of a real-photograph probe, at the cost of \gateInCost{} points of
coverage on the supported set. On \webN{} web images of the twenty dishes,
screened for near-duplicates against the training data, the system is
\webTop\% top-1 accurate and accepts \webAccept\%. Three follow-up experiments
are negative: stronger photometric augmentation, class-balanced and focal losses
move accuracy by less than the seed spread, which locates the remaining error in
intra-class variation and in the boundary of the catch-all class rather than in
the optimisation.
\end{abstract}

\begin{IEEEkeywords}
Chinese food recognition, canteen tray, transfer learning, pre-training scale,
open-set recognition, calibration, calorie estimation
\end{IEEEkeywords}